\section{Police Chase}
\label{sec:cses-police-chase}

\problemstatement{Given an undirected graph of $n$ intersections and $m$
streets, find the minimum number of streets to remove so that node $1$ is
disconnected from node $n$, and output those streets.}

\begin{patternbox}
``Fewest edges to disconnect $s$ from $t$'' is the \textbf{minimum cut},
which by the max-flow--min-cut theorem equals the \textbf{maximum flow} with
every edge given capacity $1$. The cut edges are those crossing from the
source-reachable side to the rest after the flow saturates.
\end{patternbox}

\subsection*{Max flow, then read the cut}

Model each street as an undirected unit-capacity edge (a pair of directed
edges). The max flow from $1$ to $n$ equals the minimum number of streets
whose removal disconnects them. After computing it, BFS in the residual
graph from node $1$: the reachable set $S$ is one side of the cut. Every
original street with exactly one endpoint in $S$ is a cut edge -- print
those. Their count equals the max-flow value.

\begin{ideabox}[The Reachable Side Defines the Cut]
When no augmenting path remains, the nodes still reachable from the source
in the residual graph are separated from the rest by saturated edges. The
original edges straddling that boundary are precisely a minimum edge cut.
\end{ideabox}

\subsection*{Algorithm}

\begin{enumerate}
  \item Add each street as a unit-capacity edge in both directions; run max
        flow $1\to n$.
  \item BFS the residual graph from $1$ to get the reachable set $S$.
  \item Output every original street with one endpoint in $S$ and one
        outside.
\end{enumerate}

\subsection*{Dry run}

\dryrun{Path $1\!-\!2\!-\!3$ with $n=3$; a single street between $2$ and
$3$ (and $1,2$).}

\begin{itemize}
  \item Max flow $=1$; residual-reachable set from $1$ is $\{1,2\}$.
  \item The street $2\!-\!3$ crosses the boundary -- remove it (one
        street).
\end{itemize}

\subsection*{C++ implementation}

\begin{lstlisting}
#include <bits/stdc++.h>
using namespace std;

struct Dinic {
    struct Edge { int to; long long cap; int rev; };
    vector<vector<Edge>> g;
    vector<int> level, it;
    Dinic(int n) : g(n), level(n), it(n) {}
    void addEdge(int u, int v, long long c) {
        g[u].push_back({v, c, (int)g[v].size()});
        g[v].push_back({u, c, (int)g[u].size() - 1});   // undirected: both cap c
    }
    bool bfs(int s, int t) {
        fill(level.begin(), level.end(), -1);
        queue<int> q; level[s] = 0; q.push(s);
        while (!q.empty()) {
            int u = q.front(); q.pop();
            for (auto& e : g[u])
                if (e.cap > 0 && level[e.to] < 0) { level[e.to] = level[u] + 1; q.push(e.to); }
        }
        return level[t] >= 0;
    }
    long long dfs(int u, int t, long long f) {
        if (u == t) return f;
        for (int& i = it[u]; i < (int)g[u].size(); i++) {
            Edge& e = g[u][i];
            if (e.cap > 0 && level[u] < level[e.to]) {
                long long d = dfs(e.to, t, min(f, e.cap));
                if (d > 0) { e.cap -= d; g[e.to][e.rev].cap += d; return d; }
            }
        }
        return 0;
    }
    long long maxflow(int s, int t) {
        long long flow = 0;
        while (bfs(s, t)) {
            fill(it.begin(), it.end(), 0);
            long long f;
            while ((f = dfs(s, t, LLONG_MAX)) > 0) flow += f;
        }
        return flow;
    }
};

int main() {
    int n, m;
    cin >> n >> m;
    Dinic dinic(n + 1);
    vector<pair<int,int>> streets(m);
    for (auto& [a, b] : streets) { cin >> a >> b; dinic.addEdge(a, b, 1); }

    long long cut = dinic.maxflow(1, n);
    // residual reachability from source
    vector<bool> reach(n + 1, false);
    queue<int> q; q.push(1); reach[1] = true;
    while (!q.empty()) {
        int u = q.front(); q.pop();
        for (auto& e : dinic.g[u])
            if (e.cap > 0 && !reach[e.to]) { reach[e.to] = true; q.push(e.to); }
    }

    cout << cut << "\n";
    for (auto& [a, b] : streets)
        if (reach[a] != reach[b]) cout << a << " " << b << "\n";
    return 0;
}
\end{lstlisting}

\begin{complexitybox}
\textbf{Time Complexity.} $O(V^2E)$ for the flow; the cut extraction is
$O(V+E)$. \textbf{Space Complexity.} $O(V+E)$.
\end{complexitybox}

\begin{takeawaybox}
Minimum edge cut equals maximum flow with unit capacities, and the cut edges
straddle the residual-reachable frontier. Max-flow--min-cut is one of the
most powerful reductions in graph theory -- ``disconnect cheaply'' becomes
``push maximally.''
\end{takeawaybox}
